% Physical page 5 onward: the developed commentary of the concise study,
% organized by the questions the three readings answer differently rather than
% by the order of the Mass. Each unit draws on several appointed elements and
% sets the readings' answers beside one another.

\section{The Propers: Detailed Commentary}
\zlabel{triptych:concise:commentary:start}

\subsection{Two commandments, and what holds them together}

The lawyer asks for one commandment and receives two. St Augustine takes the
threefold ``whole'' of the first as an exclusion: it means ``that no part of
our life is to be unoccupied, and to afford room, as it were, for the wish to
enjoy some other object'', so that whatever else presents
itself as worthy of love ``is to be borne into the same channel in which the
whole current of our affections flows'' (\work{De doctrina christiana}
I.22.20--21). Love of neighbour is then not a second stream beside the first;
it has an order, ``every man is to be loved as a man for God's sake; but God is
to be loved for His own sake'' (I.27.28), and a reach that takes in every man,
the angels, and the Lord himself, who in the parable of the man left half dead
``desired to be called our neighbor'' (I.30.31--33).

St Gregory the Great states the same ``whole'' as a contrast of measures.
Expounding the wedding garment of the parable that opens the same chapter of
Matthew, he calls it charity, held between two beams like woven
cloth, \latin{sicut in duobus lignis}. In the love of neighbour a
measure is given, ``as thyself''; the love of God is bound by none, for the
commandment says not how much but out of how much: \latin{quia ille veraciter
Deum diligit, qui sibi de se nihil relinquit}, he truly loves God who leaves
nothing of himself to himself (\work{Homiliae in Evangelia} 38.10). Gregory
cites the commandment in St Mark's form and expounds the parable rather than
this pericope; what he supplies is the sharpest account of the measureless
love beside Augustine's.

St John Chrysostom answers the other question the text raises, in what sense
the second commandment is ``like'' the first: ``Because this makes the way for
that, and by it is again established'' (\work{Homilies on Matthew} 71). Both
clauses speak of the second commandment: the love of neighbour opens the road
to the love of God, and the love of God in its turn re-establishes the love of
neighbour. He also notices that Jesus answers more than he was asked, and reads
the unbidden second commandment as exposing where the question came from ---
``from having no charity, from pining with envy''. St Jerome gives another
account of the motive: the question was captious because every commandment of
God is great, so that any answer could be attacked. St Thomas
Aquinas gathers these and adds
his own: the first commandment is \latin{maximum capacitate}, greatest in what
it holds, since the love of neighbour is contained in it, and the second is
like it because man is made to God's likeness, so that in the neighbour
\latin{diligitur Deus in illo}, God himself is loved (\work{Super Matthaeum}
c.~22, lect.~4).

The four agree that the commandments cannot be separated and differ only in
what they make the link: one current of love carrying everything to one good
(Augustine), a measure set against measurelessness (Gregory), mutual dependence
in both directions (Chrysostom), the image of God in the neighbour (Aquinas).
The
Collect asks as a gift what the Gospel commands, \latin{te solum Deum pura
mente sectari}; Augustine joins following to loving --- ``We
follow after God by loving Him'' --- and calls the first commandment the one
``which leads to happy life'' (\work{De moribus Ecclesiae catholicae}
I.11.18). Read from the Introit instead, the same Collect sounds different: the
people that asks to follow God alone has just confessed that God is just and
begged to be dealt with by mercy and not by its own justice, so that its love
of God begins as mercy received.

\subsection{``What think you of Christ?''}

The Gospel's second question is the one that ends the questioning, and on it
the witnesses are unanimous. St Jerome presses the logic of the psalm: if the
Christ is a mere man and only David's son, how can David call him Lord, and
that not by error nor by his own will but in the Holy Spirit? \latin{Dominus
igitur David vocatur, non secundum id quod de eo natus est, sed juxta id quod
natus ex Patre semper fuit, praeveniens ipsum carnis suae patrem} --- he is
called David's Lord by the birth from the Father that always was, and that goes
before the father of his flesh (\work{Commentarii in Matthaeum} IV, on
22:41--46).

St Augustine makes the answer a confession. Christians answer as the
questioners did, ``David's'', \latin{sed non remaneamus ubi Judaei}; the
answer they could not give is the prologue of St John, so that in the form of
God he is David's Lord and in the form of a servant David's son. Christ, he
says, \latin{non te filium ejus negasti, sed modum quo id fieret inquisisti},
did not deny the sonship but asked how it came to be; and the questioners
\latin{maluerunt inflata taciturnitate disrumpi, quam humili confessione
edoceri}, preferred to burst with a swollen silence rather than be taught by a
humble confession (\work{Enarrationes in Psalmos} 109, 3--6). Aquinas states
the doctrine: the psalm shows the Son's \latin{praeeminentiam ad sanctos,
aequalitatem ad patrem, et dominium super rebelles}, so that he is
\latin{filius secundum carnem \dots\ et dominus secundum divinitatem}.

Chrysostom supplies what joins this half of the Gospel to the other. Jesus had
just praised the confession that there is one God, and the saying ``There is
one God'' is spoken ``not to the rejection of the Son, but to make the
distinction from idols''; the counter-question then leads the questioners ``to
confess Him also to be God''. On that reading the God whom the first
commandment requires to be loved with the whole heart is not loved apart from
his Christ, and the second half of the Gospel belongs to the reading anchored
in the commandment as well. The reading anchored in the
Introit enters here only at the end, through Augustine's contrast: where the
questioners kept a proud silence, the people that prays this Mass confesses.
William Durandus, whose own Sunday kept another Gospel, takes its second half
as Christ reproving the Pharisees \latin{de infidelitate, quia non credebant ipsum esse Deum}
(\work{Rationale} VI, on the Eighteenth Sunday).

\subsection{The sevenfold \latin{unus}, and the division at its last verse}

Paul's appeal is the same charity seen from the side of a community, and its
first interpreters read verses 1--3 as a school of humility. Chrysostom begins
from the calling: ``Ye were called as His body'', and lowliness is ``the basis
of all virtue'', since one who remembers ``that all is of grace'' will humble
himself; forbearance follows from the same mercy, ``If thou \dots\ art not
forbearing to thy neighbor, how shall God be forbearing to thee?'' The unity is
God's work before it is ours: ``to this end was the Spirit given, that He might
unite those who are separated by race and by different manners''
(\work{Homilies on Ephesians} 9). St Jerome reads the same verses with an eye for the weak: true
humility is \latin{non tam in sermone, quam in mente}, ``supporting one another
in charity'' means bearing the sick brother, the poor and the widow, and the
command to keep the unity ``in the bond of peace'' tells against those who,
having broken the bond, suppose that they still hold the unity of the Spirit
(\work{Commentarii in epistulam ad Ephesios} II). Theodoret of Cyrus explains
why lowliness comes first: the gifts of the Spirit could puff up those who
received them, whereas one grace is given to all and one spring sends out the
different streams.
Aquinas names the end of the whole exhortation, that unity be kept: charity is
\latin{conjunctio animorum}, its bond is peace, and one who falls short is not
corrected at once but \latin{misericorditer expectari}, awaited with mercy.

At verse 6 the witnesses divide into two lines of exegesis, and neither denies
the other. Jerome, writing against Sabellius and the Arians, gives the three
phrases to the three Persons:
\latin{Super omnes enim est Deus Pater, quia auctor est omnium. Per omnes
Filius, quia cuncta transcurrit, vaditque per omnia. In omnibus Spiritus
sanctus, quia nihil absque eo est}. Chrysostom instead reads them of the one
God and Father: ``\,`Who is over all,' that is, the Lord and above all; and
`through all,' that is, providing for, ordering all; and `in you all,' that is,
who dwelleth in you all. Now this they own to be an attribute of the Son; so
that were it an argument of inferiority, it never would have been said of the
Father.'' Theodoret names the same three relations --- lordship, providence and
indwelling --- and draws from the repeated ``one'' of verses 5--6 that
brothers owe one another concord. Aquinas holds both: he commends the one
God's dignity \latin{ex tribus}, the height of his godhead, the breadth of his
power and the largesse of his grace, then appropriates the first phrase to the
Father as \latin{fontale principium divinitatis}, the second to the Son, the
third to the Holy Spirit.

The two readings of the Epistle take the verse accordingly. The reading of the
whole heart and the one body rests on verses 1--4 and on the moral the Greek
commentators draw from the sevenfold ``one'': Chrysostom's ``God hath called
you all on the same terms \dots\ whence is it that ye are high-minded?'' and
Theodoret's concord among brothers. The reading of David's son and David's
Lord rests on verses 5--6 for the confession of the one God, on Jerome and
Aquinas for the distribution among the Persons, and on Chrysostom for an
argument about the Son rather than for that distribution.

\subsection{The nation God chose, and the word by which the heavens stand}

The Gradual carries two verses that the three readings pull in three
directions. On the respond, verse 12, Augustine finds the happiness all men want and answers it in
two words, \latin{Hoc ama, hoc posside}: God is his people's possession and
they are his inheritance, \latin{alterutrum se possident}, each possessing the
other. Cassiodorus hears the nation as the people gathered from all nations,
an inheritance acquired and not bequeathed, \latin{quem praedicationibus
sanctis et pretioso sanguine conquisivit}, won by holy preaching and precious
blood. Aquinas adds that the beatitude which all desire, to cleave to God by
knowledge and love, is held now \latin{inchoative et in spe} and perfectly
hereafter. Read from the Introit's confession, the same verse says something
plainer: Augustine writes of this nation that \latin{neque haec a seipsa, sed
Dei munere electa est}, it was chosen not of itself but by God's gift.

The versicle, verse 6, gathers a consensus that the respond does not. St
Irenaeus of Lyons uses it to state the rule of faith against those who divided
God: there is one God Almighty ``who made all things by His Word \dots\ `By
the Word of the Lord were the heavens established, and all the might of them, by
the spirit of His mouth''' (\work{Adversus haereses} I.22.1). St Basil the
Great reads its terms directly: ``You are therefore to perceive three, the Lord
who gives the order, the Word who creates, and the Spirit who confirms''
(\work{De Spiritu Sancto} 16.38). Augustine stops his hearers at the same
verse --- \latin{Verbum certe Dei Filius est, et Spiritus oris ejus Spiritus
sanctus est} --- and at once guards the Father's part: \latin{Numquid sine
Patre?} Cassiodorus finds the Trinity in the grammar, \latin{Verbo} declaring
the Son, \latin{Domini} the Father and \latin{spiritu oris eius} the Holy
Spirit, \latin{et ut in tribus personis manifestam intelligeres unitatem, eius
dixit, non eorum}. Aquinas reads the verse by the same appropriation he uses at
Eph 4:6, and St Robert Bellarmine keeps the literal sense, God's word of
command in Genesis 1, beside the mystery: ``there is no doubt but the Holy
Ghost meant to glance at the mystery of the Holy Trinity to be revealed in the
New Testament''. The Missal prints this versicle here with a small
letter, \latin{spíritu}, and \latin{Spíritu} where it sets the same Gradual in
the votive Mass of the Holy Spirit; Basil, Augustine and Cassiodorus read the
Word and the Spirit out of the psalm's own words.

\subsection{The voice in which this people prays}

The Alleluia is one verse long, and its witnesses differ over whose voice it
is. For Augustine the poor man of Ps~101 is Christ, head and body:
\latin{Vox ergo una, quia caro una}, one voice because one flesh, a reading he
supports from the very chapter of Ephesians this Sunday reads, ``Until we all
meet into the unity of faith \dots\ unto a perfect man''; and the head speaks
for every member, \latin{quicumque in meo corpore tribulantur, ego tribulor}.
In the same place he asks \latin{unde illi labor, unde gemitus Verbo, per quod
facta sunt omnia?}, whence labour and groaning should come to the Word through
whom all things were made, and answers that the Word deigned to take our death.
In the same sentence he marks off three unions by the gender of one pronoun
--- Christ and the Church \latin{unus}, one man; the Father and the Word
\latin{unum}, one in substance; the Word and the flesh not \latin{unum} in that
sense, but one person --- and that distinction is what the reading of David's
Lord needs, since the one Christ is David's son in the flesh he took and
David's Lord in the substance he has of the Father.
Cassiodorus declines to put the psalm in Christ's mouth, since Christ is
nowhere said to have been ``anxious'', and hears instead the poor of Christ, \latin{qui non
solum pro suis malis, verum etiam pro totius mundi calamitatibus intercedere
comprobantur}, interceding not only for their own ills but for the calamities
of the whole world. Each answer serves its own reading: Augustine's one voice
of head and members serves the reading of the one body,
the Word made poor for us that of David's Lord, and Cassiodorus's interceding
poor that of the humbled people.
Augustine supplies that third reading a second sentence on the same verse
---\,\latin{si me extollo, longe fis; si me humilio, inclinas aurem tuam ad
me}\,--- and Bellarmine notes that the Church uses the verse daily as a
preparation for every other petition.

At the Offertory the disagreement touches the praying subject himself. St
Jerome observes that Daniel confesses the people's sins in
the first person \latin{quia unus e populo est}, because he is one of the
people, and marks the logic of the prayer: first \latin{Tibi, Domine,
justitia}, since \latin{juste enim patimur, quod meremur}, and then, because
the Lord is not only just but merciful, an appeal \latin{post sententiam
judicantis} to his clemency; at the petition the Offertory sings he writes
\latin{Imple opere, quod sermone pollicitus es}. Where Daniel later says ``my
sins, and the sins of my people'', Jerome allows a second possibility, that the
prophet joined himself to a sinful people out of humility though he had not
sinned. Augustine will not allow it, and rests on the words themselves:
Daniel ``expresses himself in language so distinct and precise \dots\ `My
sins,' says he, `and the sins of my people.''' Aquinas cites the verse the antiphon does not sing, ``not for our
justifications \dots\ but for the multitude of thy tender mercies'', as the
ground of the confidence every prayer needs. The reading of David's Lord takes
the same chant further forward, to the answer Gabriel brought while Daniel was
still speaking, as Rupert of Deutz and Durandus read it; but the 1962 antiphon
ends at the petitions of verses 17--19 and prints no verses, so that weight
rests on words this Mass does not sing.

The Introit stands behind both. Augustine reads verse 137 as the penitent's
reason for weeping, since God's justice \latin{omni est metuenda peccanti}, and
verse 124 as the plea that follows from it, \latin{fac cum servo tuo secundum
misericordiam tuam: non utique secundum iustitiam meam}, the
\latin{iustificationes} being those \latin{quibus Deus facit iustos, non ipsi
se}. St Hilary of Poitiers reads the same two verses for perseverance rather
than for beginning: the servant needs God's pity \latin{ut in hac seruitutis
nostrae professione maneamus}, and mercy is given \latin{ut uolentes adiuuet,
incipientes confirmet, adeuntes recipiat}. The two differ where the beginning
lies: Hilary, writing before the controversy with the Pelagians, says
\latin{ex nobis autem initium est, ut illa perficiat}, where Augustine gives the
whole of it to God; they agree that the Introit asks mercy and not reward. Cassiodorus gives the verse the same use and adds a link of his own, that by this saying Daniel was
saved among the lions --- which joins the Introit to the prophet whose prayer
the Offertory sings.

\subsection{Vows paid, vices cured, and the end this Mass looks to}

The Communion turns the day's confession into an act. Augustine counts first
the vows common to all Christians, to believe, to hope for eternal life and to
live well, before the particular vows some freely add, and reads ``all you that
are round about him'' of those who hold the truth in common and humbly;
\latin{Spiritus enim principum, superbi sunt spiritus}, and the kings of the
earth he takes morally, \latin{Rege terram, et eris rex terrae}. Cassiodorus
and Bellarmine hear the same words where the chant is sung, of the faithful who
bring their gifts to the altar. The reading of David's Lord hears the antiphon
against the Gospel's own psalm instead: the God who takes away the spirit of
princes is the Lord at whose right hand the Son sits until every enemy is made
his footstool, and Aquinas's \latin{dominium super rebelles} is the phrase for
it.

In all three readings the three orations are prayer rather than the anchor of
an argument, and two of the readings ask what the sacrament
those prayers seek actually does. For that both call on one chapter of the
Council of Trent (sess.~XIII, cap.~2): the reading of the one body for the
Eucharist as \latin{symbolum unius illius corporis} whose head is Christ, and
the reading of the humbled people for the \latin{antidotum} that frees from
daily faults. Nothing in the reading of David's Lord
rests on that doctrine. That reading hears two of the three prayers
as confession: the Collect's \latin{te solum Deum}, which Chrysostom's point
keeps from being read against the Son, and the Secret's \latin{maiestas}, which
meets the Preface the Missal directs immediately after it, confessing ``one
God, and one Lord not in a singularity of one Person, but in a Trinity of one
substance''. The Postcommunion adds nothing distinctive to it: the address is
to the \latin{omnípotens Deus} with no
distinction of Persons, and what it asks is the healing the other two readings
develop.

Where the Mass looks at the end is where the readings differ most quietly.
Jerome's ``one hope'' of the calling is the kingdom, the one house of the
Father with its many mansions, and Bellarmine reads the blessed nation as most happy
``when we shall see him as he is''. Rupert of Deutz, expounding the Introit as
his own books gave it, hears its petition asking \latin{hanc videlicet
misericordiam, ut dicas illi: Amice, ascende superius}, the mercy by which the
host at the banquet says ``Friend, go up higher''; Durandus hears its
\latin{rectum iudicium} as the judgment \latin{quo Deus exaltat humiles, et
superbos deprimit}. Both read these chants within an office whose Gospel
was Luke~14. Bellarmine takes the Introit's \latin{iudicium} otherwise, holding
that throughout this psalm ``the word `judgment' always means the law''; they
differ in what the word names, not in what they believe of God. He reads the
verses just before the Communion's, ``When God arose in judgment, to save all
the meek of the earth'', of the general
judgment, and the Postcommunion asks that the remedies it seeks be eternal.

All three readings end at the same place, the silence with which the Gospel
closes. ``No man was able to answer him a word.'' The first answers that
silence with love, the second with confession, the third with prayer; and the
Mass itself, having heard the Gospel, answers it with the Creed.
